% Target: rmp-iii-a-pentagon-mpo
% Formula: ^1F ^1F = ^0F ^1F ^1F — recoupling the two fusion trees.
% Ink: Kernel flat frame; two directed three-leaf fusion trees.
% Boundary: Each tree keeps one incoming root and three outgoing leaves open.
% Source: paper TN-Review-main.tex line 1317; paper figure fig3_pentagon2.pdf;
%   no author drawing source
% Capabilities: kernel, coherence, recoupling
\[
  \tenkzkernel
  \tndeclare{species}{fus}{hue=source:red}
  \tndeclare{species}{vertex}{hue=source:lightgray}
  \begin{tenkzeq}
    \begin{tenkz}[rows={wire,wire,wire,wire}, cols=4, bonds=none]
      % the left tree fuses the upper pair first
      \tn[at=(3,1), name=A, skin=dot, species=vertex,
          ports={180:virtual, 45:virtual, 315:virtual}]{}
      \tn[at=(2,2), name=B, skin=dot, species=vertex,
          ports={225:virtual, 45:virtual, 0:virtual}]{}
      \tnwire[species=fus, dir=to]{open w}{A.180}
      \tnwire[species=fus, dir=to]{A.45}{B.225}
      \tnwire[species=fus, dir=to]{B.45}{(1,4)}
      \tnwire[species=fus, dir=to]{B.0}{(2,4)}
      \tnwire[species=fus, dir=to]{A.315}{(4,4)}
    \end{tenkz}
    =
    \begin{tenkz}[rows={wire,wire,wire,wire}, cols=4, bonds=none]
      % the right tree fuses the lower pair first
      \tn[at=(2,1), name=C, skin=dot, species=vertex,
          ports={180:virtual, 45:virtual, 315:virtual}]{}
      \tn[at=(3,2), name=D, skin=dot, species=vertex,
          ports={135:virtual, 45:virtual, 0:virtual}]{}
      \tnwire[species=fus, dir=to]{open w}{C.180}
      \tnwire[species=fus, dir=to]{C.45}{(1,4)}
      \tnwire[species=fus, dir=to]{C.315}{D.135}
      \tnwire[species=fus, dir=to]{D.45}{(2,4)}
      \tnwire[species=fus, dir=to]{D.0}{(3,4)}
    \end{tenkz}
  \end{tenkzeq}
\]
\[
  {}^{1}F\,{}^{1}F \;=\; {}^{0}F\,{}^{1}F\,{}^{1}F
\]
